% =====================================================================
%  CARNET D'ENTRAÎNEMENT — CHIMIE — FICHE 3 — THÈME 1
%  Particules subatomiques — NIVEAU FACILE — ÉNONCÉS
% =====================================================================
\documentclass[11pt,a4paper]{article}

\usepackage[utf8]{inputenc}
\usepackage[T1]{fontenc}
\usepackage{lmodern}
\usepackage[french]{babel}

\usepackage{amsmath,amssymb,mathtools}
\usepackage{xcolor}
\usepackage{colortbl}
\usepackage{tikz}
\usetikzlibrary{arrows.meta,positioning,calc}
\usepackage[most]{tcolorbox}
\usepackage{enumitem}
\usepackage{booktabs}
\usepackage{array}
\usepackage[a4paper,margin=2.1cm,headheight=26pt,headsep=9pt]{geometry}
\usepackage{fancyhdr}
\usepackage[hidelinks]{hyperref}

\definecolor{primary}{RGB}{124,78,140}
\definecolor{primarydark}{RGB}{74,44,92}
\definecolor{defcol}{RGB}{0,114,178}
\definecolor{thmcol}{RGB}{0,140,120}
\definecolor{attncol}{RGB}{200,40,55}

\newcommand{\nuc}[3]{\prescript{#1}{#2}{\mathrm{#3}}}
\newcommand{\pp}{p^{+}}
\newcommand{\ee}{e^{-}}
\newcommand{\nz}{n^{0}}

\pagestyle{fancy}
\fancyhf{}
\fancyhead[L]{\small\textcolor{primarydark}{\textbf{Fiche 3 — Thème 1}} : Particules subatomiques}
\fancyhead[R]{\small\textcolor{primarydark}{\textsc{Facile — Énoncés}}}
\fancyfoot[C]{\small\thepage}
\renewcommand{\headrulewidth}{0.8pt}
\renewcommand{\headrule}{\hbox to\headwidth{\color{primary}\leaders\hrule height \headrulewidth\hfill}}

\newcounter{exo}
\newcommand{\exo}[1][]{\stepcounter{exo}\par\medskip%
  \noindent{\color{primary}\rule[-2pt]{3.5pt}{15pt}}\ \textbf{\color{primarydark}\large Exercice \theexo}%
  \ifx\relax#1\relax\relax\else\ \textmd{\normalsize\color{black!65}— \itshape #1}\fi\par\nopagebreak\smallskip}

\setlist{leftmargin=1.7em,topsep=2pt,itemsep=3pt}

\begin{document}

\begin{tikzpicture}
\node[rectangle,rounded corners=6pt,fill=primary,text=white,
      minimum width=\textwidth,minimum height=1.35cm,align=center,inner sep=6pt]
  {\Large\bfseries Thème 1 — Particules subatomiques\\[2pt]
   \normalsize\mdseries Niveau Facile — Énoncés};
\end{tikzpicture}

\vspace{3pt}
{\small\itshape Chaque exercice isole une seule mécanique. Réponds sans calculatrice quand c'est possible.
Données : $e = 1{,}6\times10^{-19}$ C ; $m(\pp)\approx m(\nz)\approx 1{,}7\times10^{-27}$ kg ;
$m(\ee)\approx 9{,}1\times10^{-31}$ kg.}

% ---------------------------------------------------------------
\exo[carte d'identité des trois particules]
Pour chacune des trois particules — \textbf{proton}, \textbf{neutron}, \textbf{électron} — indique :
\begin{enumerate}[label=\alph*)]
  \item sa \textbf{localisation} dans l'atome (noyau ou nuage électronique) ;
  \item le \textbf{signe} de sa charge (positive, négative ou nulle).
\end{enumerate}

% ---------------------------------------------------------------
\exo[du plus lourd au plus léger]
Range les trois particules (proton, neutron, électron) \textbf{de la plus lourde à la plus légère}.
Deux d'entre elles ont une masse quasiment identique : lesquelles ?

% ---------------------------------------------------------------
\exo[vrai ou faux]
Indique si chaque affirmation est \textbf{vraie} ou \textbf{fausse}. Justifie chaque réponse en une phrase.
\begin{enumerate}[label=\alph*)]
  \item Le neutron porte une charge négative.
  \item L'électron n'a pas de masse.
  \item Le proton et l'électron portent des charges de signes opposés mais de même valeur absolue.
  \item Un électron est environ deux mille fois plus léger qu'un proton.
\end{enumerate}

% ---------------------------------------------------------------
\exo[charges élémentaires]
Donne, en coulombs :
\begin{enumerate}[label=\alph*)]
  \item la charge d'un seul proton ;
  \item la charge d'un seul électron ;
  \item la charge totale d'un ensemble de \textbf{5 protons}.
\end{enumerate}

% ---------------------------------------------------------------
\exo[le constituant principal d'un atome]
Le diamètre d'un noyau est de l'ordre de $10^{-15}$ m ; celui de l'atome entier de l'ordre de
$10^{-10}$ m. Parmi les protons, les neutrons, les électrons et le vide, quel est le \textbf{constituant
principal, en volume,} d'un atome quelconque ? Justifie ta réponse en une phrase.

\end{document}
